% fig102-fenwquery.tex -- Figure 5.x: Fenwick prefix(7) reads nodes 7, 6, 4.
% fstyle.tex -- shared TikZ styles for the figures of this book (palette from lifeos.sty).
\tikzset{%
  cell/.style={draw=ink, minimum width=7mm, minimum height=6.5mm, inner sep=0pt, font=\footnotesize, fill=white},
  cellh/.style={cell, fill=accentsoft, draw=accent},
  cella/.style={cell, fill=amberbg, draw=amber},
  cellg/.style={cell, fill=goodbg, draw=good},
  cellx/.style={cell, fill=graybg, draw=slate, text=slate},
  vtx/.style={circle, draw=ink, minimum size=6.5mm, inner sep=0pt, font=\footnotesize, fill=white},
  vtxh/.style={vtx, draw=accent, fill=accentsoft},
  vtxa/.style={vtx, draw=amber, fill=amberbg},
  vtxg/.style={vtx, draw=good, fill=goodbg},
  vtxd/.style={vtx, draw=slate, fill=graybg, text=slate},
  lab/.style={font=\footnotesize, text=slate},
  labA/.style={font=\footnotesize, text=accent},
  labM/.style={font=\footnotesize, text=amber},
  idx/.style={font=\tiny, text=slate},
  arr/.style={-{stealth}, thick, ink},
  arrA/.style={-{stealth}, thick, accent},
  arrM/.style={-{stealth}, thick, amber},
  arrG/.style={-{stealth}, thick, good},
  edge/.style={thick, ink},
  edgeA/.style={thick, accent},
  edgeM/.style={thick, amber},
}%
% \cellrow{prefix}{x0}{y}{style}{v0,v1,...}: a row of cells named prefix0, prefix1, ... (pitch 0.75cm)
\providecommand{\cellrow}[5]{\foreach \v [count=\i from 0] in {#5} {\node[#4] (#1\i) at (#2+\i*0.75,#3) {\v};}}%
% \cellidx{prefix}{n}{start}: index labels under cells prefix0 .. prefix(n-1)
\providecommand{\cellidx}[3]{\foreach \i in {0,...,\numexpr#2-1\relax} {\node[idx, below=0pt of #1\i] {\the\numexpr\i+#3\relax};}}%
% \fitfig{<tikzpicture>}: shrink the picture to \linewidth only when it would overflow.
\ifdefined\fitbox\else\newsavebox{\fitbox}\fi%
\providecommand{\fitfig}[1]{\sbox{\fitbox}{#1}\ifdim\wd\fitbox>\linewidth\resizebox{\linewidth}{!}{\usebox{\fitbox}}\else\usebox{\fitbox}\fi}%

\fitfig{\begin{tikzpicture}[font=\footnotesize, x=1.1cm, y=0.78cm]
\foreach \i/\lb in {1/1,2/2,3/1,4/4,5/1,6/2,7/1,8/8} {
  \pgfmathsetmacro{\lev}{ln(\lb)/ln(2)}
  \pgfmathsetmacro{\lo}{\i-\lb}
  \ifnum\i=7 \def\fillc{amber!25}\def\drc{amber}\else\ifnum\i=6 \def\fillc{amber!25}\def\drc{amber}\else\ifnum\i=4 \def\fillc{amber!25}\def\drc{amber}\else\def\fillc{graybg}\def\drc{slate}\fi\fi\fi
  \draw[draw=\drc, fill=\fillc] (\lo+0.06,\lev+0.1) rectangle (\i-0.06,\lev+0.9);
  \node[font=\scriptsize] at ({(\lo+\i)/2},\lev+0.5) {$t[\i]$};}
\foreach \i in {1,...,8} \node[idx, below] at (\i-0.5,0) {\i};
\node[labM, anchor=west, align=left] at (8.4,2.6) {prefix$(7)$: $i=7$, read $t[7]$, $i\leftarrow 7-1=6$;\\ read $t[6]$, $i\leftarrow6-2=4$; read $t[4]$, $i\leftarrow 0$};
\node[lab, anchor=west, align=left] at (8.4,0.9) {node $i$ covers the last $\mathrm{lowbit}(i)$ cells;\\ strip the lowest set bit each step};
\end{tikzpicture}}
